% ===== Chapter 2, The gift of memory (~40 min; part 2 of webinar 1 until the 2026-10-23 split) =====
% 33 pt instead of the macro's 38 (2026-11-12): at 38 the one-word title ran into the mascot circle
\partopener{Գլուխ 2}{\texorpdfstring{\fontsize{33}{35}\selectfont Հիշողություն}{Հիշողություն}}{Հիշողության նվերը (պահեք կտրոնը)}{thinking}{lavender}

% DRM word list: shown now, tested at the start of "when memory fails"
\begin{frame}{Հիշողության փոքրիկ թեստ}
  \begin{columns}[c]
    \begin{column}{0.3\textwidth}
      \centering
      \mascot[2.6cm]{thinking}\par\medskip
      Կարդա ու փորձիր հիշել։\\ \textbf{Մի գրիր}։\\ Կվերադառնանք դրան։
    \end{column}
    \begin{column}{0.64\textwidth}
      \begin{card}
        \large
        \begin{tabular}{@{}lll@{}}
          մահճակալ & մրափ & խաղաղություն\\[1.4mm]
          հանգիստ & վերմակ & հորանջ\\[1.4mm]
          արթուն & նինջ & թմրած\\[1.4mm]
          հոգնած & խռմփոց & գիշեր\\[1.4mm]
          երազ & բարձ & արթնանալ
        \end{tabular}
      \end{card}
    \end{column}
  \end{columns}
\end{frame}

% Reworded 2026-11-12 (Hayk: no Care Bears, clearer); the bottom line says what the chapter will show
\begin{frame}{Հիշողությունը համակարգիչ չէ}
  \begin{columns}[c]
    \begin{column}{0.6\textwidth}
      \begin{popup}{հիշողություն.exe}
        Պատկերացրու համակարգիչ, որը...
        \begin{enumerate}
          \item<2-> \textbf{ինքն է որոշում}՝ ինչն է կարևոր
          \item<3-> ֆայլերը \textbf{խառնում է} իրար
          \item<4-> ամոթալի ֆայլերդ \textbf{ինքն է բացում}
          \item<5-> չհավանածը \textbf{փոխում է} իր ճաշակով
        \end{enumerate}
      \end{popup}
    \end{column}
    \begin{column}{0.36\textwidth}
      \onslide<6->{%
      \begin{card}[colback=brainpink!30!cardwhite]
        Նման համակարգիչը կես ժամից կնետեիր պատուհանից \emoji{computer}\par\smallskip
        Իսկ ուղեղդ \textbf{հենց այսպես է} աշխատում, ու ոչ փոխել ես կարող, ոչ անջատել-միացնել։
      \end{card}}
    \end{column}
  \end{columns}
  \vspace{1mm}
  \onslide<7->{\centering\small Այսօր կտեսնենք՝ ինչու է ուղեղն այսպես անում, ու ինչու է այն, միևնույն է, աշխատում։\par}
\end{frame}

% ---------------------------------------------------------------- short-term vs long-term
\topicopener{door}{Կարճ ու երկար}{Ինչի՞ համար էի եկել այստեղ}{Կարճաժամկետ և երկարաժամկետ հիշողություն}{confused}

\begin{frame}{Ինչի՞ համար էի եկել}
  \begin{columns}[c]
    \begin{column}{0.48\textwidth}
      \begin{poll}
        Բարձրացրու ձեռքդ, եթե վերջին շաբաթվա ընթացքում մտել ես սենյակ ու մոռացել, թե ինչու։
      \end{poll}
    \end{column}
    \begin{column}{0.48\textwidth}
      \onslide<2->{%
      \begin{tikzpicture}[storystep/.style={flow=5.4cm, align=left, inner sep=4pt}]
        \node[storystep] (a) at (0,0) {\emoji{bulb}\ «Գնամ խոհանոցից մկրատը բերեմ»};
        \node[storystep, visible on=<3->] (b) at (0,-1.55) {\emoji{musical}\ Ճանապարհին՝ սիրածդ երգը, մեմ, սերիալի սյուժեի լուծումը};
        \node[storystep, visible on=<4->, fill=brainpink!30!cardwhite] (c) at (0,-3.1) {\emoji{door}\ Խոհանոց. «...ինչի՞ համար էի եկել»};
        \draw[-{Stealth}, plum, line width=1pt, visible on=<3->] (a) -- (b);
        \draw[-{Stealth}, plum, line width=1pt, visible on=<4->] (b) -- (c);
      \end{tikzpicture}}
    \end{column}
  \end{columns}
  \vspace{2mm}
  \onslide<5->{%
  \begin{card}[colback=lavender!35!cardwhite]
    \centering Նոր միտքը \textbf{դուրս է հրել} հինը կարճաժամկետ հիշողությունից, մինչ այն կհասցներ «պահպանվել»։
  \end{card}}
\end{frame}

\begin{frame}{Երկու տարբեր հիշողություն}
  \begin{columns}[T]
    \begin{column}{0.47\textwidth}
      \begin{card}[colback=peach!35!cardwhite]
        {\blackfont\large Կարճաժամկետ} \emoji{coffee}\par
        short-term / working memory\par
        \begin{itemize}
          \item Առանց կրկնելու՝ \textbf{մինչև 1 րոպե}
          \item \textbf{Մոտ 4} միավոր
          \item Նեյրոնների \textbf{ակտիվություն}. ցուցակ կապուչինոյի փրփուրի վրա՝ մի պահ կա, հետո չկա
          \item Հենց սրանով ենք \textbf{մտածում}
        \end{itemize}
      \end{card}
    \end{column}
    \begin{column}{0.47\textwidth}
      \onslide<2->{%
      \begin{card}[colback=mint!35!cardwhite]
        {\blackfont\large Երկարաժամկետ} \emoji{books}\par
        long-term memory\par
        \begin{itemize}
          \item Կարող է մնալ \textbf{ամբողջ կյանք}
          \item Սահմանն անհայտ է. ոչ ոք դեռ չի լցրել այն
          \item Նոր \textbf{կապեր} նեյրոնների միջև (սինապսներ)
          \item Մտածելու «նյութ» է
        \end{itemize}
      \end{card}}
    \end{column}
  \end{columns}
  \vspace{1mm}
  \onslide<3->{\centering Երեկվա դասախոսությունն արդեն \textbf{երկարաժամկետ} հիշողություն է։\par}
\end{frame}

\begin{frame}{4, ոչ թե 7}
  \begin{columns}[T]
    \begin{column}{0.46\textwidth}
      \begin{itemize}
        \item 1950-ականներին Ջորջ Միլլերը գրեց «կախարդական թվի» մասին. 7 ± 2 միավոր (Miller's law)
        \item Ավելի ճշգրիտ փորձերը ցույց տվեցին՝ \textbf{մոտ 4}
        \item Բայց «միավորը» կարող է մեծ լինել. ուղեղը խմբավորում է (chunking)
      \end{itemize}
    \end{column}
    \begin{column}{0.48\textwidth}
      \begin{tryit}
        Հիշիր այս թվերը.\par\smallskip
        {\blackfont\Large 0 9 3 1 8 2 5 0 9}\par\smallskip
        \onslide<2->{Իսկ հիմա՝\par\smallskip {\blackfont\Large 093 \quad 182 \quad 509}\par\smallskip
        Նույն 9 թվանշանն են, բայց հիմա ընդամենը 3 միավոր։}
      \end{tryit}
    \end{column}
  \end{columns}
  \source{Miller, 1956; Cowan, 2010}
\end{frame}

% Our example since 2026-11-12 (Hayk: not family-related; the book's is "your mum smells of cheese"),
% plus the point of chunking and the book's translator example, so the slide says why it matters
\begin{frame}{Chunking. ինչու է «միավորը» կարևոր}
  \centering
  \chip{}{բլոգեր}\quad\chip{}{կատուդ}\quad\chip{}{է}\quad\chip{}{TikTok}\quad\chip{}{գաղտնի}\quad{\small $\rightarrow$ 5 միավոր}\par\smallskip
  \onslide<2->{%
  \chip[butter]{blackcat}{Կատուդ գաղտնի TikTok բլոգեր է}\quad{\small $\rightarrow$ 1 միավոր}\par\smallskip
  {\small ...ու մեկ կասկածելի հայացք կատվիդ կողմը}\par\smallskip}
  \onslide<3->{%
  \begin{card}[width=0.92\textwidth, colback=lavender!35!cardwhite]
    \raggedright\textbf{Իմաստը}. կարճաժամկետը փոքր է, բայց \textbf{երկարաժամկետը} նրան պատրաստի «բլոկներ» է տալիս՝ բառեր, արտահայտություններ, թվեր։ Ինչքան շատ ես գիտես, այնքան մեծ են բլոկները։
  \end{card}}
  \onslide<4->{%
  \begin{card}[width=0.92\textwidth, colback=mint!35!cardwhite]
    \raggedright\emoji{speech}\ Սինխրոն թարգմանիչն ամեն բառ առանձին չի պահում. բառերն ու քերականությունը երկարաժամկետում են, իսկ կարճաժամկետը միայն հերթականությունն ու իմաստն է «շարում»։
  \end{card}}
\end{frame}

% Transition (2026-11-12, Hayk: no abrupt topic changes)
\begin{frame}
  \bigquestion{Լավ, իսկ ֆիզիկապես որտե՞ղ է «գրված» հիշողությունը}{Նեյրոններ, սինապսներ ու հիպոկամպ}
\end{frame}

\begin{frame}{Որտեղ է «գրված» հիշողությունը}
  \centering
  \begin{tikzpicture}[lbl/.style={font=\small, text=plum, align=center}]
    % neuron 1: dendrites, cell body, axon, terminal
    \foreach \a in {120,160,200,240}{\draw[plum, line width=1.2pt] (-5.2,0) -- ++(\a:1.0) -- ++(\a+25:0.4);}
    \filldraw[fill=brainpink, draw=plum, line width=1.2pt] (-5.2,0) circle (0.55);
    \draw[plum, line width=2.2pt] (-4.65,0) -- (-1.05,0);
    \filldraw[fill=brainpink, draw=plum, line width=1.2pt] (-0.75,0) ellipse (0.35 and 0.55);
    % labels kept clear of the dendrites (2026-11-12): "neuron" right of the lowest one, the signal label on two lines
    \node[lbl] at (-4.6,-1.0) {նեյրոն};
    \draw[-{Stealth}, amOrange!85!black, line width=1.4pt] (-4.3,0.35) -- node[above, lbl] {էլեկտրական ազդանշան\\ (action potential)} (-1.5,0.35);
    % synapse gap with neurotransmitters
    \foreach \y/\x in {0.3/-0.2, 0/0.0, -0.25/-0.15, 0.15/0.15, -0.4/0.1}{\fill[amBlue, visible on=<2->] (\x,\y) circle (0.06);}
    % neuron 2 with receptors
    \filldraw[fill=lavender, draw=plum, line width=1.2pt] (0.75,0) ellipse (0.35 and 0.55);
    \foreach \y in {-0.3,0,0.3}{\draw[amBlue, line width=1.2pt, visible on=<2->] (0.43,\y-0.08) -- (0.36,\y) -- (0.43,\y+0.08);}
    \draw[plum, line width=2.2pt] (1.1,0) -- (4.6,0);
    \filldraw[fill=lavender, draw=plum, line width=1.2pt] (5.2,0) circle (0.55);
    \foreach \a in {-60,-20,20,60}{\draw[plum, line width=1.2pt] (5.2,0) -- ++(\a:1.0) -- ++(\a-25:0.4);}
    \node[lbl, visible on=<2->] (syn) at (0,-1.45) {\textbf{սինապս}՝ նեղ ճեղք\\ նեյրոմեդիատորներ \emoji{key} ու ընկալիչներ \emoji{lock}};
    \draw[plum, line width=0.8pt, visible on=<2->] (syn.north) -- (0,-0.6);
  \end{tikzpicture}
  \begin{columns}[T]
    \begin{column}{0.47\textwidth}
      \onslide<3->{%
      \begin{card}[colback=peach!35!cardwhite]
        \textbf{Կարճաժամկետ}՝ արդեն եղած կապերի ժամանակավոր \textbf{ակտիվություն}։ Արագ, բայց անկայուն
      \end{card}}
    \end{column}
    \begin{column}{0.47\textwidth}
      \onslide<4->{%
      \begin{card}[colback=mint!35!cardwhite]
        \textbf{Երկարաժամկետ}՝ \textbf{նոր սինապսներ} (encoding)։ Ուղեղը հատուկ կապեր է կառուցում, ու դա ժամանակ է պահանջում
      \end{card}}
    \end{column}
  \end{columns}
\end{frame}

\begin{frame}{Հիպոկամպը՝ հիշողությունների գործարան}
  \centering
  \begin{tikzpicture}[
      st/.style={flow=2.75cm, line width=0.9pt, rounded corners=2.5mm, inner sep=4pt, minimum height=1.25cm},
      ar/.style={-{Stealth[length=3mm]}, plum, line width=1.4pt},
      arlbl/.style={font=\scriptsize, text=plum, align=center, fill=cream, inner sep=1pt},
    ]
    \node[st] (sen) at (-5.8,0) {\emoji{eye}\emoji{ear}\\ \textbf{Զգայարաններ}};
    \node[st, fill=peach!40!cardwhite] (stm) at (-2.0,0) {\emoji{coffee}\\ \textbf{Կարճաժամկետ}\\ մտածում ենք};
    \node[st, fill=brainpink!40!cardwhite, visible on=<2->] (hip) at (1.8,0) {\emoji{brain}\\ \textbf{Հիպոկամպ}\\ կոդավորում};
    \node[st, fill=mint!40!cardwhite, visible on=<3->] (ctx) at (5.6,0) {\emoji{books}\\ \textbf{Կեղև}\\ երկար պահպանում};
    \draw[ar] (sen) -- (stm);
    \draw[ar, visible on=<2->] (stm) -- (hip);
    \node[arlbl, visible on=<2->] at (-0.1,1.35) {կրկնություն \emoji{cyclone}\\ զգացմունք \emoji{heart}};
    \draw[ar, visible on=<3->] (hip) -- (ctx);
    \node[arlbl, visible on=<3->] at (3.7,1.35) {կամաց «տեղափոխվում» է\\ (consolidation)};
    \draw[ar, dashed, plum!60] (stm.south) -- ++(0,-0.75) node[below, font=\scriptsize, text=plum!80] {չկրկնվածը՝ թույլ է պահվում}; % was «մնացածը մոռացվում է», which contradicted the next frame (2026-11-12)
  \end{tikzpicture}
  \begin{columns}[T]
    \begin{column}{0.47\textwidth}
      \onslide<4->{%
      \begin{card}\emoji{taxi}\ Լոնդոնի տաքսու վարորդների \textbf{հետին հիպոկամպն} է մեծ\end{card}}
    \end{column}
    \begin{column}{0.47\textwidth}
      \onslide<5->{%
      \begin{card}\emoji{pin}\ Համարը կրկնելը «Շտապ» պիտակ է. «Սա պահիր առանձին»\end{card}}
    \end{column}
  \end{columns}
\end{frame}

% Transition (2026-11-12): the book's own question at this point
\begin{frame}
  \bigquestion{Եթե երկարաժամկետը գրեթե ամեն ինչ պահում է, ինչո՞ւ ենք մոռանում}{«Լավ հարց է»՝ Բըրնեթ}
\end{frame}

\begin{frame}{Ինչու ենք մոռանում}
  \begin{columns}[c]
    \begin{column}{0.52\textwidth}
      {\blackfont Հիշողության 3 քայլը}\par\smallskip
      \chip[peach]{pencil}{1. Գրանցել}\ \chip[mint]{books}{2. Պահել}\ \chip[lavender]{magnifier}{3. Գտնել}\par\medskip
      encoding $\rightarrow$ storage $\rightarrow$ retrieval\par\medskip
      \onslide<2->{Մոռացածի մեծ մասը, ամենայն հավանականությամբ, \textbf{դեռ այնտեղ է}։ Պարզապես չենք գտնում։}
    \end{column}
    \begin{column}{0.42\textwidth}
      \onslide<3->{%
      \begin{bookquote}
        Ինչպես ձեռնոցները, որ չես գտնում. ձեռնոցներ ունես, բայց ձեռքերդ միևնույն է մրսում են։
      \end{bookquote}
      \centering\bigemoji[1.2cm]{gloves}}
    \end{column}
  \end{columns}
\end{frame}

\begin{frame}{Ինչը հեշտ է գտնվում, ինչը՝ ոչ}
  \begin{columns}[T]
    \begin{column}{0.47\textwidth}
      \begin{card}[colback=mint!35!cardwhite]
        {\blackfont Հեշտ է}\par\smallskip
        Ինչը կապված է զգացմունքի հետ (շատ կապեր, շատ «ճանապարհներ»)։ Առաջին համբույրը, կամ երբ ավտոմատը մեկի գնով \textbf{երկու} չիփս տվեց \emoji{chips}
      \end{card}
    \end{column}
    \begin{column}{0.47\textwidth}
      \onslide<2->{%
      \begin{card}[colback=lavender!35!cardwhite]
        {\blackfont Դժվար է}\par\smallskip
        Ինչը ոչ մի բանի հետ կապված չէ։\par\smallskip
        Դասի գնալու 509-րդ ձանձրալի ճանապարհը
      \end{card}}
    \end{column}
  \end{columns}
  \vspace{2mm}
  \onslide<3->{%
  \begin{card}[colback=peach!35!cardwhite]
    Ծայրահեղ դեպք՝ տրավմայի հիշողություններ, որ նորից ու նորից են «վերադառնում» (գրքում՝ «flashbulb memories»). այնքան վառ են, որ խանգարում են ապրել։
  \end{card}}
\end{frame}

\begin{frame}{Համատեքստը հիշողության բանալին է}
  \begin{columns}[c]
    \begin{column}{0.52\textwidth}
      1975. սուզորդները բառեր են սովորել \textbf{ցամաքում} կամ \textbf{ջրի տակ}, ու ստուգվել նույն կամ մյուս տեղում\par\smallskip
      \centering
      \begin{tikzpicture}[cell/.style={draw=plum, line width=0.8pt, minimum width=2.1cm, minimum height=0.9cm,
          rounded corners=1.5mm, font=\small, align=center}]
        \node[font=\blackfont\scriptsize, align=center] at (0,1.05) {ստուգում՝\\ ցամաքում \emoji{house}};
        \node[font=\blackfont\scriptsize, align=center] at (2.3,1.05) {ստուգում՝\\ ջրի տակ \emoji{diving}};
        \node[font=\blackfont\scriptsize, align=center, anchor=east] at (-1.2,0) {սովորել՝\\ ցամաքում};
        \node[font=\blackfont\scriptsize, align=center, anchor=east] at (-1.2,-1.15) {սովորել՝\\ ջրի տակ};
        \node[cell, fill=mint!50!cardwhite] at (0,0) {\emoji{check}\ ավելի լավ};
        \node[cell, fill=brainpink!30!cardwhite] at (2.3,0) {ավելի վատ};
        \node[cell, fill=brainpink!30!cardwhite] at (0,-1.15) {ավելի վատ};
        \node[cell, fill=mint!50!cardwhite] at (2.3,-1.15) {\emoji{check}\ ավելի լավ};
      \end{tikzpicture}
    \end{column}
    \begin{column}{0.42\textwidth}
      \begin{itemize}
        \item Նույն միջավայրը «միացնում է» հիշողության մի մասը (context-dependent memory)
        \item Ինչպես կախաղան խաղի բացված տառերը
      \end{itemize}
      \onslide<2->{%
      \begin{scienceupdate}
        2021-ին 16 սուզորդով կրկնված փորձը \textbf{չհաստատեց} արդյունքը։ Դասական փորձերն էլ են վերստուգվում։
      \end{scienceupdate}}
    \end{column}
  \end{columns}
  \source{Godden \& Baddeley, 1975; Murre, 2021}
\end{frame}

\begin{frame}{Հիշողության տեսակները}
  \centering
  \begin{tikzpicture}[kind/.style={flow=3.8cm, line width=0.9pt, rounded corners=2.5mm, inner sep=4pt, anchor=north}]
    \node[kind, fill=butter!50!cardwhite, text width=4.4cm] (lt) at (0,3.0) {\textbf{Երկարաժամկետ հիշողություն}};
    \node[kind, fill=peach!40!cardwhite, visible on=<2->] (epi) at (-4.5,1.6) {\emoji[1.4em]{camera}\ \textbf{Էպիզոդիկ}\\ (episodic) քեզ հետ կատարվածը\\ \emph{«Էյֆելյան աշտարակից ներքև փսխեցի»}};
    \node[kind, fill=mint!40!cardwhite, visible on=<3->] (sem) at (0,1.6) {\emoji[1.4em]{books}\ \textbf{Իմաստային}\\ (semantic) փաստեր՝ առանց համատեքստի\\ \emph{«Փարիզը Ֆրանսիայի մայրաքաղաքն է»}};
    \node[kind, fill=lavender!45!cardwhite, visible on=<4->] (pro) at (4.5,1.6) {\emoji[1.4em]{videogame}\ \textbf{Պրոցեդուրային}\\ (procedural) հմտություններ\\ \emph{հեծանիվ քշել, տպել առանց նայելու}};
    \draw[plum, line width=1pt, visible on=<2->] (lt.south) -- (epi.north);
    \draw[plum, line width=1pt, visible on=<3->] (lt.south) -- (sem.north);
    \draw[plum, line width=1pt, visible on=<4->] (lt.south) -- (pro.north);
  \end{tikzpicture}
  \vspace{0.5mm}
  \onslide<5->{%
  \begin{card}[colback=lavender!35!cardwhite]
    \centering Բըրնեթը պրոցեդուրային հիշողության մասին չի խորանում. «եթե սկսեք դրա մասին մտածել, կարող է դժվարանաք օգտագործել այն»
  \end{card}}
\end{frame}

% ---------------------------------------------------------------- faces and names
\topicopener{eyes}{Դեմքեր ու անուններ}{Էյ, դու... էն... էն անգամվանից}{Ինչու դեմքը հիշում ենք, իսկ անունը՝ ոչ}{thinking}

\begin{frame}{Գիտե՞ս էն աղջկան...}
  \begin{columns}[c]
    \begin{column}{0.72\textwidth}
      \begin{tikzpicture}
        \node[chatmsg=8cm, fill=peach!40!cardwhite, anchor=north west] (m1) at (0,0)
          {\textbf{Մայրիկ.} Հիշո՞ւմ ես էն բարձրահասակ աղջկան՝ քո դասարանից։ Մազերը ներկում էր, մեր կողքի փողոցում էին ապրում, մինչև ծնողները բաժանվեցին։ Քույրը քո զարմիկի ընկերուհին էր։ Միշտ կարմիր վերարկու էր հագնում, չէր սազում իրեն...};
        \node[chatmsg=2.4cm, fill=cardwhite, anchor=north east, visible on=<2->] (m2) at ([yshift=-3mm]m1.south east) {\textbf{Ես.} Անո՞ւնը};
        \node[chatmsg=2.4cm, fill=peach!40!cardwhite, anchor=north west, visible on=<3->] at ([yshift=-3mm]m1.south west |- m2.south) {\textbf{Մայրիկ.} Չգիտեմ։};
      \end{tikzpicture}
    \end{column}
    \begin{column}{0.24\textwidth}
      \centering
      \onslide<3->{\mascot[2.4cm]{confused}}
    \end{column}
  \end{columns}
\end{frame}

\begin{frame}{Մայրիկների հիշողության կանոնը}
  \centering
  \begin{meme}[6.2cm]{drake.jpg}
    \memecap[text width=2.6cm]{0.75}{0.75}{Հիշել անունը}
    \memecap[text width=2.7cm, font=\blackfont\footnotesize]{0.75}{0.25}{Հիշել մազերի գույնը, կարմիր վերարկուն, քրոջ ընկերներին ու ծնողների բաժանվելը}
  \end{meme}
\end{frame}

\begin{frame}{Դեմքը նկար է, անունը՝ ձայն}
  \begin{columns}[T]
    \begin{column}{0.47\textwidth}
      \begin{card}[colback=mint!35!cardwhite]
        \emoji[1.3em]{eyes}\ {\blackfont Դեմք}\par\smallskip
        \begin{itemize}
          \item Աչքեր, մազեր, դիմախաղ. շատ տեղեկություն
          \item Ուղեղն ունի դեմքերի հատուկ «սարքավորում»
          \item Երկարաժամկետ հիշողությունը սիրում է \textbf{պատկերներ ու իմաստ}
        \end{itemize}
      \end{card}
    \end{column}
    \begin{column}{0.47\textwidth}
      \onslide<2->{%
      \begin{card}[colback=brainpink!25!cardwhite]
        \emoji[1.3em]{ear}\ {\blackfont Անուն}\par\smallskip
        \begin{itemize}
          \item Պատահական վանկեր, որ կպցված են դեմքին
          \item Ապրում է \textbf{կարճաժամկետում}, պետք է կրկնել
          \item Իսկ զրույցը («որտեղի՞ց ես», «ի՞նչ ես սովորում») նրան դուրս է հրում
        \end{itemize}
      \end{card}}
    \end{column}
  \end{columns}
  \vspace{0.5mm}
  \onslide<3->{\centering\small Անունը սովորաբար հնչում է \textbf{առաջինը}, ու ուղեղը լավ է հիշում ցուցակի առաջին ու վերջին բաները (primacy, recency)։ Սա մի քիչ օգնում է։\par}
\end{frame}

\begin{frame}{«Այս ֆայլն արդեն գոյություն ունի»}
  \begin{columns}[c]
    \begin{column}{0.5\textwidth}
      \begin{popup}{դեմքեր.exe}
        {\blackfont Այս դեմքն արդեն գոյություն ունի}\par\smallskip
        Որտեղի՞ց, ե՞րբ, ի՞նչ անուն. տվյալներ չկան։
      \end{popup}
    \end{column}
    \begin{column}{0.44\textwidth}
      \begin{itemize}
        \item<2-> \textbf{Ճանաչում} (familiarity)՝ գիտես, որ հանդիպել ես, ու վերջ
        \item<3-> \textbf{Վերհիշում} (recall)՝ գտնում ես ամբողջ հիշողությունը՝ ով է, որտեղից, անունը
        \item<4-> Ճանաչումն էժան է ու արագ. ծանոթը նշանակում էր «ինձ չսպանեց», կարելի է կենտրոնանալ նորի վրա
      \end{itemize}
    \end{column}
  \end{columns}
\end{frame}

\begin{frame}{«Լեզվիս ծայրին է»}
  \begin{columns}[c]
    \begin{column}{0.46\textwidth}
      \centering
      \begin{tikzpicture}
        % wider bar pitch and a two-line label (2026-11-12): the four labels ran into each other at 1.15 cm
        \draw[plum, line width=0.8pt] (0,0) -- (5.6,0);
        \foreach \i/\h/\t in {1/1.0/դեմք, 2/1.6/{+ որտեղ}, 3/2.2/{+ ում\\ հետ}, 4/3.0/{+ «Բա...»}}{
          \fill[amBlue, visible on=<\i->] (\i*1.3-0.85,0) rectangle ++(0.75,\h);
          \node[font=\scriptsize, text=plum, below, align=center, visible on=<\i->] at (\i*1.3-0.475,0) {\t};}
        \draw[amRed, dashed, line width=1.2pt] (0,2.6) -- (5.6,2.6) node[right, font=\scriptsize, text=amRed] {շեմ};
        \node[visible on=<5->, font=\blackfont, text=amRed] at (4.7,3.45) {«Բաբկեն»};
      \end{tikzpicture}
    \end{column}
    \begin{column}{0.48\textwidth}
      \begin{itemize}
        \item Ամեն հուշում մի քիչ է «տաքացնում» հիշողությունը
        \item Վերհիշումը միանում է միայն \textbf{շեմն} անցնելուց հետո (recall threshold)
      \end{itemize}
      \onslide<5->{%
      \begin{bookquote}
        Հարևանի հրավառությունը լուսավորում է մութ տունը, մինչև ներսիններն արթնանում են ու վառում բոլոր լույսերը։
      \end{bookquote}}
    \end{column}
  \end{columns}
\end{frame}

\begin{frame}{Ինչպես հիշել անունը}
  \begin{columns}[c]
    \begin{column}{0.58\textwidth}
      \begin{enumerate}
        \item Լսելուն պես \textbf{կրկնիր}. «Ուրախ եմ, Բաբկեն»
        \item Մի քանի անգամ \textbf{օգտագործիր} զրույցում
        \item Կապիր \textbf{պատկերի} հետ, որ անունն էլ դառնա «նկար»
      \end{enumerate}
      \onslide<2->{%
      \begin{bookquote}
        Եթե նորից հանդիպենք ու անունդ չհիշեմ, կոպիտ չեմ։ Լավ, կոպիտ եմ։ Բայց հիմա գիտես՝ ինչու։
      \end{bookquote}}
    \end{column}
    \begin{column}{0.36\textwidth}
      \centering
      \mascot[2.8cm]{happy}
    \end{column}
  \end{columns}
\end{frame}

% ---------------------------------------------------------------- questions pause
% Was the break slide; no break since 2026-11-12 (Hayk), he asks for questions here. Glazz has letters only: no ՞ in the title
\begin{frame}[plain]
  \begin{tikzpicture}[remember picture, overlay]
    \fill[butter!60!cream] (current page.south west) rectangle (current page.north east);
    \node at ([yshift=8mm]current page.center) {\mascot[4cm]{thinking}};
    \node[anchor=north, font=\titlefont\fontsize{30}{32}\selectfont, text=plum] at ([yshift=-17mm]current page.center) {Դադար};
    \node[anchor=north, font=\mediumfont\large, text=plum] at ([yshift=-31mm]current page.center)
      {Հարցե՞ր կան մինչև հիմա ասածի մասին \emoji{speech}};
    \node[anchor=north, font=\mediumfont, text=plum!70!cream] at ([yshift=-39mm]current page.center)
      {Գրեք չաթում կամ միացրեք խոսափողը};
  \end{tikzpicture}
\end{frame}

% ---------------------------------------------------------------- alcohol
\topicopener{wine}{Ալկոհոլ}{Մի բաժակ գինի՝ հիշողությունը թարմացնելու համար}{Ինչպես է ալկոհոլը ջնջում հիշողությունը, և երբ է տարօրինակ կերպով օգնում}{sick}

\begin{frame}{Ալկոհոլն անջատում է «կարմիր լույսերը»}
  \begin{columns}[c]
    \begin{column}{0.56\textwidth}
      % 2026-11-12: the paradox spelled out (Hayk did not get the leg quote's point)
      \begin{itemize}
        \item Էթանոլը փոքր է ու լուծվում է ջրում. հեշտությամբ հասնում է ուղեղ
        \item Այն \textbf{դեպրեսանտ} է. ճնշում է նեյրոնների ակտիվությունը, ինչպես ձայնն իջեցնելը \emoji{mute}
        \item<2-> Ուրեմն խմածը պետք է հանգիստ նստի՞։ Ոչ. ուղեղի մեծ մասը զբաղված է \textbf{արգելակելով} (inhibition), ինչպես կարմիր լույսերը մեծ քաղաքում \emoji{trafficlight}
        \item<4-> Ալկոհոլը նախ թուլացնում է հենց այդ \textbf{«արգելակները»}, ու մարդն անում է այն, ինչ սովորաբար զսպում էր
      \end{itemize}
    \end{column}
    \begin{column}{0.38\textwidth}
      \onslide<3->{%
      {\blackfont\small Արգելակման օրինակ}\par\smallskip
      \begin{bookquote}
        Ոտքդ շարժող մասը կարևոր է, բայց ոչ հանդիպման ժամանակ, երբ պետք է հանգիստ նստես։ Դրա համար մեկ այլ մաս նրան արգելակում է. «Հիմա պետք չէ, ախպեր»։
      \end{bookquote}}
    \end{column}
  \end{columns}
\end{frame}

\begin{frame}{Ինչն է անջատվում առաջինը}
  \begin{columns}[c]
    \begin{column}{0.5\textwidth}
      \begin{tikzpicture}[lev/.style={flow=5.6cm, align=left, minimum height=0.68cm}]
        \node[lev, fill=brainpink!45!cardwhite] at (0,0) {\textbf{1.} զսպվածություն, ամոթ, «սա լավ միտք չէ»};
        \node[lev, fill=peach!45!cardwhite, visible on=<2->] at (0,-0.85) {\textbf{2.} խոսք, քայլվածք};
        \node[lev, fill=butter!60!cardwhite, visible on=<3->] at (0,-1.7) {\textbf{3.} հիշողության գրանցում};
        \node[lev, fill=amRed!25!cardwhite, visible on=<4->] at (0,-2.55) {\textbf{4.} շնչառություն, սիրտ. \textbf{կյանքին վտանգ}};
      \end{tikzpicture}
    \end{column}
    \begin{column}{0.44\textwidth}
      \begin{card}[colback=butter!45!cream]
        Ինչքան \textbf{նոր ու բարդ} է համակարգը, այնքան շուտ է «ընկնում»։
      \end{card}
      \onslide<5->{%
      \begin{bookquote}
        80-ականների Walkman-ը կդիմանա աստիճաններից ընկնելուն, իսկ սմարթֆոնը կջարդվի սեղանի անկյունից։
      \end{bookquote}}
    \end{column}
  \end{columns}
\end{frame}

% Split in two on 2026-11-12 (Hayk: «խոռոչ», not «անցք»; what «անջատվել» means; the video-game point):
% this frame defines gaps vs a full blackout, the next one is the book's alcoholics with the game analogy
\begin{frame}{Հիշողության «սև խոռոչներ»}
  \begin{columns}[c]
    \begin{column}{0.52\textwidth}
      \begin{itemize}
        \item Ալկոհոլը հատկապես խանգարում է \textbf{հիպոկամպին}՝ նոր հիշողությունների «գործարանին»
        \item<2-> Հաճախ հիշողությունները դեռ գրանցվում են, բայց վատ ու պատահական. առավոտյան՝ «խոռոչներ»
        \item<3-> Ավելի շատ խմելիս՝ \textbf{լրիվ բլեքաութ} (blackout). մարդն արթուն է, խոսում ու քայլում է, բայց ոչինչ չի գրանցվում։ Հաջորդ օրն այդ ժամերը \textbf{չկան}
      \end{itemize}
    \end{column}
    \begin{column}{0.44\textwidth}
      \centering
      \begin{tikzpicture}[lbl/.style={font=\small, text=plum, anchor=south west, inner sep=1pt}]
        \node[lbl, font=\blackfont\small] at (0,2.75) {Երեկոն, ինչպես հիշում ես առավոտյան};
        \node[lbl, visible on=<2->] at (0,1.9) {«Խոռոչներ»};
        \fill[mint, visible on=<2->] (0,1.4) rectangle (5.6,1.85);
        \foreach \a/\b in {1.0/1.5, 2.5/2.8, 3.7/4.5}{\fill[plum, visible on=<2->] (\a,1.4) rectangle (\b,1.85);}
        \node[lbl, visible on=<3->] at (0,0.75) {Լրիվ բլեքաութ};
        \fill[mint, visible on=<3->] (0,0.25) rectangle (5.6,0.7);
        \fill[plum, visible on=<3->] (1.8,0.25) rectangle (5.6,0.7);
        \node[font=\scriptsize, text=plum!75, anchor=north west, visible on=<2->] at (0,0.1) {կանաչ՝ հիշում ես, մուգ՝ դատարկ};
      \end{tikzpicture}
    \end{column}
  \end{columns}
\end{frame}

\begin{frame}{Բլեքաութ, որ դրսից չի երևում}
  \begin{columns}[c]
    \begin{column}{0.5\textwidth}
      \begin{itemize}
        \item Գրքի խոսքով՝ մարդկանց մեծ մասը լրիվ բլեքաութի հասնելու համար այնքան պետք է խմի, որ հազիվ կանգնի
        \item<2-> Ալկոհոլից կախվածություն ունեցողների ուղեղը \textbf{հարմարվել} է. նրանք նորմալ խոսում ու շարժվում են
        \item<3-> Բայց հիշողություն չի գրանցվում. 10 րոպե անց չեն հիշում՝ ինչ ասացին կամ արեցին
      \end{itemize}
    \end{column}
    \begin{column}{0.46\textwidth}
      \onslide<4->{%
      \begin{card}[colback=lavender!35!cardwhite]
        \centering\bigemoji[1cm]{videogame}\par\smallskip
        \raggedright Խաղում ես համակարգչային խաղ, գնում ես զուգարան, ու ընկերդ շարունակում է քո փոխարեն։ Կողքից նայողը տարբերություն չի տեսնում, բայց վերադառնալիս \textbf{չգիտես՝ ինչ եղավ}։\par\smallskip
        Բլեքաութի ժամանակ էլ մարմինը «խաղում» է, իսկ հիշողությունը «դուրս է եկել»։
      \end{card}}
    \end{column}
  \end{columns}
\end{frame}

% Myth or fact (2026-11-12): the book implies only alcoholics black out while still talking; students do too
\begin{frame}{Առասպել, թե՞ փաստ}
  % statement = the book's own claim (frame before), since 2026-11-12 review: "during a blackout you can barely stand"
  \mythorfact{Բլեքաութի ժամանակ մարդը հազիվ է կանգնում}{ԱՌԱՍՊԵԼ}{pinkdark}{%
    % one line only: longer text ran off the frame (the statement box and stamp leave no room); the rest is in NOTES
    \vspace{-4mm}
    \begin{scienceupdate}[fontupper=\footnotesize]
      Մարդը խոսում ու քայլում է։ Խմող ուսանողների \textbf{կեսն} այն ունեցել է
    \end{scienceupdate}}
  \source{White et al., J. Am. Coll. Health, 2002; NIAAA}
\end{frame}

\begin{frame}{Երբ «վիճակն» է օգնում հիշել}
  \begin{columns}[T]
    \begin{column}{0.54\textwidth}
      \begin{itemize}
        \item Ոչ միայն միջավայրը, այլև \textbf{ներքին վիճակն} է համատեքստ. ալկոհոլ, կոֆեին, տրամադրություն (state-dependent recall)
        \item<2-> Վատ տրամադրությամբ սովորածն ավելի հեշտ ես հիշում նորից վատ տրամադրությամբ
        \item<3-> Երկու բաժակ գինուց հետո լսած բամբասանքն ավելի լավ կհիշես մեկ այլ երեկո, նորից երկու բաժակից հետո։ Այդ դեպքում գինին իսկապես «թարմացնում է հիշողությունը»
      \end{itemize}
    \end{column}
    \begin{column}{0.4\textwidth}
      \onslide<4->{%
      \begin{bookquote}
        Խնդրում եմ, սա մի ընդունեք որպես գիտական խորհուրդ քննությունից առաջ խմելու համար։ Հատկապես եթե վարորդական իրավունքի քննություն է։
      \end{bookquote}}
    \end{column}
  \end{columns}
\end{frame}

\begin{frame}{Բըրնեթի քննությունը}
  \begin{columns}[c]
    \begin{column}{0.56\textwidth}
      \begin{itemize}
        \item Ամբողջ գիշեր սովորեց սուրճով, ու քննությունից առաջ էլ մի մեծ բաժակ խմեց \emoji{coffee}
        \item<2-> Արդյունքը՝ \textbf{73\%}, կուրսի լավագույններից մեկը
        \item<3-> Բայց ամբողջ քննության ընթացքում զուգարան էր ուզում
        \item<4-> Լրացուցիչ թուղթ խնդրելիս քննողին ասաց \textbf{«պապա»}
        \item<5-> Իսկ տուն գնալիս կատաղի վիճեց... \textbf{աղավնու} հետ \emoji{dove}
      \end{itemize}
    \end{column}
    \begin{column}{0.38\textwidth}
      \centering
      \mascot[2.8cm]{panic}\par\medskip
      \onslide<5->{{\small Ինքն էլ է խորհուրդ տալիս \textbf{չկրկնել}}}
    \end{column}
  \end{columns}
\end{frame}

% ---------------------------------------------------------------- ego bias
\topicopener{crown}{Եսասեր հիշողություն}{Իհարկե հիշում եմ, իմ գաղափարն էր}{Ինչպես է հիշողությունը քեզ մի քիչ ավելի լավը ցույց տալիս}{smug}

\begin{frame}{Ուղեղդ դու ես}
  \begin{columns}[c]
    \begin{column}{0.58\textwidth}
      \begin{itemize}
        \item Հիշողությունը գրքի էջ չէ. այն \textbf{ճկուն} է (plastic) ու անընդհատ խմբագրվում է
        \item Ուղեղդ քեզ համար աստղի քծնող օգնականն է, որ բոլոր վատ լուրերը հեռու է պահում
        \item<2-> \textbf{Եսակենտրոն շեղում} (egocentric bias). խմբային որոշման մեջ յուրաքանչյուրը հիշում է, որ իր դերն ավելի մեծ էր, քան իրականում
        \item<3-> Սա \textbf{սուտ չէ}. մարդն անկեղծորեն հավատում է իր տարբերակին
      \end{itemize}
    \end{column}
    \begin{column}{0.36\textwidth}
      % hand-raise first (2026-11-12): the egocentric bias, live
      \begin{poll}
        Վերջին խմբային նախագծում \textbf{դու} արել ես քո բաժնից ավելի՞ն։ Բարձրացրու ձեռքդ։
      \end{poll}
      \onslide<2->{%
      \begin{card}[colback=brainpink!25!cardwhite]
        \centering Բոլորը չեն կարող իրենց բաժնից ավելին անել \emoji{thinking}
      \end{card}}
    \end{column}
  \end{columns}
\end{frame}

% Veritasium (Hayk asked, 2026-11-12): "Why Being Delusional is a Superpower" (2020), its first minute;
% numbers checked against the video's own transcript. Answers the poll on the previous frame.
\begin{frame}{Ո՞վ արեց աշխատանքի մեծ մասը}
  \begin{columns}[T]
    \begin{column}{0.58\textwidth}
      \begin{card}[colback=mint!30!cardwhite]
        \textbf{Համահեղինակներ}. «Քանի՞ տոկոսն է քոնը»\par\smallskip
        \centering
        \begin{tikzpicture}
          \fill[amBlue] (0,0) rectangle (3.4,0.42);
          \fill[amRed] (3.4,0) rectangle (4.76,0.42);
          \draw[plum, dashed, line width=0.9pt] (3.4,-0.15) -- (3.4,0.58) node[above, font=\scriptsize, text=plum, inner sep=1pt] {100\%};
          \node[right, font=\blackfont, text=amRed] at (4.76,0.21) {գումարը՝ 140\%};
        \end{tikzpicture}
      \end{card}
      \onslide<2->{%
      \begin{card}[colback=lavender!35!cardwhite]
        \textbf{Զույգեր}, տնային գործը. համարյա միշտ \textbf{100\%-ից ավելի}
      \end{card}}
      \onslide<3->{%
      \begin{card}[colback=brainpink!25!cardwhite]
        Պարծենալ չէ. «ո՞վ է սկսում վեճերը», «ո՞ւմ թափթփածն է»՝ նորից \textbf{100\%-ից ավելի}
      \end{card}}
    \end{column}
    \begin{column}{0.36\textwidth}
      \onslide<4->{%
      \begin{card}[colback=butter!45!cream]
        Քո արածը հիշում ես \textbf{ամբողջությամբ ու վառ}, ուրիշներինը՝ միայն տեսածդ մասը։ Դրա համար քո բաժինը մեծ է թվում՝ ու լավում, ու վատում։
      \end{card}
      \centering\href{https://youtu.be/3LopI4YeC4I}{\qrcode[height=1.05cm]{https://youtu.be/3LopI4YeC4I}}\par
      {\scriptsize Veritasium-ի վիդեոն}}
    \end{column}
  \end{columns}
  \source{Veritasium, 2020; Ross \& Sicoly, 1979; Caruso et al., 2006}
\end{frame}

\begin{frame}{Ջոն Դինի հիշողությունը}
  \begin{columns}[c]
    \begin{column}{0.56\textwidth}
      \begin{itemize}
        \item 1973, Ուոթերգեյթ. Նիքսոնի խորհրդական Ջոն Դինը մանրամասն պատմեց գաղտնի հանդիպումների մասին
        \item<2-> Հետո գտնվեցին այդ հանդիպումների \textbf{ձայնագրությունները} \emoji{musical}
        \item<3-> Ընդհանուր իմաստը ճիշտ էր, մանրամասները՝ հաճախ ոչ։ Ու այդ զրույցներում իր դերը հիշում էր \textbf{ավելի կենտրոնական}, քան էր
      \end{itemize}
    \end{column}
    \begin{column}{0.38\textwidth}
      \centering
      \bigemoji[1.4cm]{judge}\par\medskip
      \onslide<3->{%
      \begin{card}[colback=butter!45!cream]
        Չէր ստում։ Նրա հիշողությունն էր «հարմարվել» իր մասին իր պատկերացմանը։
      \end{card}}
    \end{column}
  \end{columns}
  \source{Neisser, Cognition, 1981}
\end{frame}

\begin{frame}{Շեղումներ, որ քեզ գովում են}
  \begin{columns}[T]
    \begin{column}{0.31\textwidth}
      \begin{card}[colback=mint!30!cardwhite]
        \centering\bigemoji[0.9cm]{trophy}\par\smallskip
        {\blackfont «Իմ ընտրածը լավագույնն էր»}\par
        \raggedright choice-supportive bias. քո ընտրությունը հիշում ես ավելի լավը, քան էր
      \end{card}
    \end{column}
    \begin{column}{0.31\textwidth}
      \onslide<2->{%
      \begin{card}[colback=lavender!35!cardwhite]
        \centering\bigemoji[0.9cm]{speech}\par\smallskip
        {\blackfont «Ես էի ասել»}\par
        \raggedright self-generation effect. քո ասածն ավելի լավ ես հիշում, քան ուրիշներինը
      \end{card}}
    \end{column}
    \begin{column}{0.31\textwidth}
      \onslide<3->{%
      \begin{card}[colback=peach!35!cardwhite]
        \centering\bigemoji[0.9cm]{monocle}\par\smallskip
        {\blackfont «Ես գիտեի»}\par
        \raggedright hindsight bias. անցյալը հիշում ես ավելի կանխատեսելի, քան այն թվում էր այդ պահին
      \end{card}}
    \end{column}
  \end{columns}
\end{frame}

\begin{frame}{...ու մի քանիսը ևս}
  \begin{columns}[T]
    \begin{column}{0.31\textwidth}
      \begin{card}[colback=mint!30!cardwhite]
        \centering\bigemoji[0.75cm]{people}\par\smallskip
        {\blackfont «Սեփական ռասայի» էֆեկտ}\par
        \raggedright own-race bias. այլ ռասայի դեմքերն ավելի դժվար ես տարբերում։ Հիմնական պատճառը՝ \textbf{փորձը}
      \end{card}
    \end{column}
    \begin{column}{0.31\textwidth}
      \onslide<2->{%
      \begin{card}[colback=lavender!35!cardwhite]
        \centering\bigemoji[0.75cm]{sun}\par\smallskip
        {\blackfont Վատը մարում է}\par
        \raggedright fading affect bias. վատ հիշողությունների ցավն ավելի արագ է մարում, քան լավերի ուրախությունը
      \end{card}}
    \end{column}
    \begin{column}{0.31\textwidth}
      \onslide<3->{%
      \begin{card}[colback=brainpink!25!cardwhite]
        \centering\bigemoji[0.75cm]{facepalm}\par\smallskip
        {\blackfont Բայց նաև հակառակը}\par
        \raggedright persistence. պարի հրավիրեցիր, նա ծիծաղեց, դու փախար ու ընկար տորթերի մեջ։ Ուղեղը սա հիշեցնում է տարիներ անց
      \end{card}}
    \end{column}
  \end{columns}
\end{frame}

\begin{frame}{Ինչո՞ւ է ուղեղն այդպես անում}
  \begin{columns}[c]
    \begin{column}{0.52\textwidth}
      \begin{itemize}
        \item Որոշումներ կայացնելու համար պետք է \textbf{վստահություն}. ամեն անգամ ամեն ինչ վերլուծելը շատ ժամանակ կխլեր
        \item Բոլոր հիշողությունները գրանցվում են \textbf{քո} տեսանկյունից
        \item Ինքնագնահատականը կարևոր է նորմալ գործելու համար
      \end{itemize}
    \end{column}
    \begin{column}{0.42\textwidth}
      \onslide<2->{%
      \begin{bookquote}
        «Դեռահաս ժամանակ վարել չսովորեցի, որովհետև \textbf{ծույլ} էի։ Իսկ հիմա չեմ սովորում, որովհետև \textbf{զբաղված} եմ»։
      \end{bookquote}
      Հին եսին քննադատելը հեշտ է. դա ներկան ավելի լավն է ցույց տալիս։}
    \end{column}
  \end{columns}
\end{frame}

\begin{frame}{Ամենամեծ ձուկը}
  \begin{columns}[c]
    \begin{column}{0.56\textwidth}
      \begin{itemize}
        \item Ամեն անգամ, երբ պատմում ես մի դեպք ու մի քիչ զարդարում, հիշողությունը \textbf{թարմացվում է} այդ նոր տարբերակով
        \item<2-> Հաջորդ պատմելիս՝ նորից մի քիչ, ու այսպես շարունակ
        \item<3-> Ձկնորսության ժամանակ բռնածդ երրորդ ամենամեծ ձուկը կամաց-կամաց դառնում է ամենամեծը \emoji{fish}
      \end{itemize}
    \end{column}
    \begin{column}{0.38\textwidth}
      \onslide<4->{%
      \begin{bookquote}
        Լավ կողմն այն է, որ եթե այս բաժինը լավ չհասկացաք, միևնույն է կհիշեք, որ հասկացել էիք։
      \end{bookquote}}
    \end{column}
  \end{columns}
\end{frame}

% ---------------------------------------------------------------- when memory fails
\topicopener{magnifier}{Երբ հիշողությունը փչանում է}{Որտե՞ղ եմ ես... Ո՞վ եմ ես}{Կեղծ հիշողություններ, ամնեզիա և հիվանդություններ}{scared}

% DRM reveal
\begin{frame}{Հիշողության թեստ. պատասխանը}
  \begin{columns}[c]
    \begin{column}{0.45\textwidth}
      % three words, not one (2026-11-12): asking only about «քուն» gave the trick away
      \begin{chatbox}
        Կայի՞ն ցուցակում. \textbf{«բարձ»}, \textbf{«քուն»}, \textbf{«մեքենա»}։ Ամեն մեկին գրիր \textbf{այո} կամ \textbf{ոչ}։
      \end{chatbox}
      \vspace{1mm}
      \onslide<3->{%
      \begin{card}[colback=brainpink!25!cardwhite]
        «Բարձ»՝ կար, «մեքենա»՝ չկար։ «Քուն»՝ \textbf{չկար}, բայց շատերը վստահ «հիշում» են այն։ Ուղեղն ինքն է \textbf{ավելացրել} ամենահավանական բառը։
      \end{card}}
    \end{column}
    \begin{column}{0.51\textwidth}
      \onslide<2->{%
      \begin{card}
        \begin{tabular}{@{}lll@{}}
          մահճակալ & մրափ & խաղաղություն\\[1mm]
          հանգիստ & վերմակ & հորանջ\\[1mm]
          արթուն & նինջ & թմրած\\[1mm]
          հոգնած & խռմփոց & գիշեր\\[1mm]
          երազ & բարձ & արթնանալ
        \end{tabular}
      \end{card}}
    \end{column}
  \end{columns}
  \source{Roediger \& McDermott, 1995}
\end{frame}

\begin{frame}{Մեկ բառ և նոր հիշողություն}
  \begin{columns}[T]
    \begin{column}{0.56\textwidth}
      Էլիզաբեթ Լոֆթուս, 1974. մարդիկ դիտել են ավտովթարի տեսագրություն ու պատասխանել. «Ի՞նչ արագությամբ էին գնում մեքենաները, երբ \textbf{\_\_\_\_} իրար»\par\medskip
      \begin{tikzpicture}
        \begin{axis}[
            xbar, width=6.6cm, height=4.4cm, bar width=0.4cm, bar shift=0pt, % two \addplot: without this the bars sit off their labels
            xmin=0, xmax=100, enlarge y limits=0.15, % xmax leaves room for the "66 կմ/ժ" labels inside the column
            symbolic y coords={hpvecin, xpvecin, dem, baxvecin, jaxjaxvecin},
            ytick={hpvecin, xpvecin, dem, baxvecin, jaxjaxvecin}, % not ytick=data: two \addplot
            yticklabels={հպվեցին, խփվեցին, դեմ առան, բախվեցին, ջախջախվեցին},
            tick label style={font=\small, text=plum}, axis line style={plum},
            axis x line*=bottom, axis y line*=left, xtick=\empty,
            nodes near coords={\pgfmathprintnumber[assume math mode=true]{\pgfplotspointmeta}\,կմ/ժ}, % text font, not serif math digits
            nodes near coords style={font=\blackfont\small, text=plum, anchor=west},
            point meta=x, clip=false,
          ]
          \addplot[fill=amBlue, draw=none] coordinates {(51,hpvecin) (55,xpvecin) (61,dem) (63,baxvecin)};
          \addplot[fill=amRed, draw=none] coordinates {(66,jaxjaxvecin)};
        \end{axis}
      \end{tikzpicture}\par
      {\scriptsize Միջին գնահատված արագություն (մղոն/ժամից վերածված)}
    \end{column}
    \begin{column}{0.4\textwidth}
      \onslide<2->{%
      \begin{card}[colback=brainpink!25!cardwhite]
        Մեկ շաբաթ անց հարցրել են. «Կոտրված ապակի տեսա՞ք»։ \textbf{Ապակի չկար}։\par\smallskip
        «Այո» պատասխանեցին՝\par
        «ջախջախվեցին» խմբից՝ \textbf{32\%}\par
        «խփվեցին» խմբից՝ \textbf{14\%}
      \end{card}}
    \end{column}
  \end{columns}
  \source{Loftus \& Palmer, 1974}
\end{frame}

\begin{frame}{Դատարանում}
  \begin{columns}[T]
    \begin{column}{0.48\textwidth}
      \begin{card}[colback=butter!45!cream]
        \emoji[1.3em]{cheese}\ {\blackfont Մեծ չեդեր-կողոպուտը}\par\smallskip
        «Ամբաստանյալը խանութի \textbf{մո՞տ} էր»։\par
        \onslide<2->{Թե՞ «Խանութի \textbf{որ մասում} էր նա»։ Երկրորդից հետո վկան «հիշում է». «Գորգոնզոլայի կողքին էր»։}
      \end{card}
    \end{column}
    \begin{column}{0.46\textwidth}
      \onslide<3->{%
      \begin{card}[colback=brainpink!25!cardwhite]
        \emoji[1.3em]{judge}\ {\blackfont Նադին Քուլ}\par\smallskip
        Թերապիայից հետո «հիշեց» սատանայական աղանդ, որը երբեք չէր եղել։ Դատով \textbf{միլիոններ} ստացավ։
      \end{card}}
    \end{column}
  \end{columns}
  \vspace{1mm}
  \onslide<4->{%
  \begin{card}[borderline west={1.4mm}{0pt}{accent}, fontupper=\mediumfont\small]
    \emoji{book}\ Բըրնեթ. «Ինձ խնդրեցին վկայել, որ բոլոր վկաները կարող են կեղծ հիշողություններ ունենալ։ Հրաժարվեցի. վախեցա, որ կքանդեմ արդարադատությունը»։
  \end{card}}
\end{frame}

% Transition (2026-11-12): the book's "* * *" break, from false memories in healthy brains to damaged ones
\begin{frame}
  \bigquestion{Մինչև հիմա ուղեղն առողջ էր, ու միևնույն է սխալվում էր։ Իսկ եթե վնասվի հենց ուղե՞ղը}{Վիրահատություն, վնասվածք, հիվանդություն}
\end{frame}

\begin{frame}{Պացիենտ Հ.Մ.}
  \begin{columns}[c]
    \begin{column}{0.58\textwidth}
      \begin{itemize}
        \item Հենրի Մոլեյսոն. ծանր էպիլեպսիա
        \item 1953-ին, 27 տարեկանում, հեռացրին քունքային բլթերի մի մասը՝ \textbf{երկու հիպոկամպերի} մոտ կեսի հետ
        \item<2-> Նոպաները դադարեցին։ Դադարեց նաև \textbf{նոր հիշողությունների} ստեղծումը
        \item<3-> Զրուցում էր նորմալ, բայց մի քանի րոպեից ամեն ինչ մոռանում էր
        \item<4-> Մինչև 2008-ը տասնամյակներ շարունակ օգնեց գիտնականներին հասկանալ հիշողությունը
      \end{itemize}
    \end{column}
    \begin{column}{0.36\textwidth}
      \centering\mascot[2.8cm]{confused}
    \end{column}
  \end{columns}
\end{frame}

\begin{frame}{Հայելու փորձը}
  \begin{columns}[c]
    \begin{column}{0.4\textwidth}
      \centering
      \begin{tikzpicture}[scale=0.9]
        \draw[plum, line width=1.8pt] (90:1.6) -- (234:1.6) -- (18:1.6) -- (162:1.6) -- (306:1.6) -- cycle;
        \draw[amRed, line width=0.9pt, decorate, decoration={random steps, segment length=1.5mm, amplitude=0.6mm}]
          (90:1.6) -- (234:1.6) -- (18:1.6);
        \node[font=\Large] at (2.3,1.4) {\emoji[1.4em]{mirror}};
      \end{tikzpicture}
    \end{column}
    \begin{column}{0.54\textwidth}
      \begin{itemize}
        \item Հ.Մ.-ն աստղ էր գծում՝ նայելով միայն \textbf{հայելուն}
        \item<2-> Ամեն օր \textbf{ավելի լավ} էր ստացվում
        \item<3-> Բայց ամեն անգամ վստահ էր, որ \textbf{առաջին անգամն} է անում
      \end{itemize}
      \vspace{2mm}
      \onslide<4->{%
      \begin{card}[colback=mint!30!cardwhite]
        Հմտությունների (procedural) հիշողությունն աշխատում է ուղեղի \textbf{այլ} մասերում, ու այն մնացել էր։
      \end{card}}
    \end{column}
  \end{columns}
\end{frame}

% Transition (2026-11-12, Hayk's example: from H.M. to the amnesia types)
\begin{frame}
  \bigquestion{Հ.Մ.-ը չէր կարողանում նորը գրանցել։ Իսկ կարելի՞ է կորցնել հինը}{Ամնեզիայի երկու տեսակ}
\end{frame}

\begin{frame}{Ամնեզիայի երկու տեսակ}
  \centering
  % labels are filled nodes on two lines since 2026-11-12: the fixed rectangles were narrower than the text
  \begin{tikzpicture}[amn/.style={font=\blackfont\small, align=center, rounded corners=1.5mm, inner sep=4pt,
      minimum width=4.6cm, anchor=south}]
    \draw[-{Stealth[length=3mm]}, plum, line width=1.6pt] (-6.4,0) -- (5.6,0) node[right, font=\small] {ժամանակ};
    \node at (0,0.05) {\emoji[1.6em]{cyclone}};
    \node[font=\blackfont\small, text=plum, below=3mm] at (0,0) {վնասվածք};
    \node[amn, fill=lavender, text=lavenderdark, visible on=<2->] at (-3.4,0.3) {Ռետրոգրադ\\ կորցնում ես \textbf{անցյալը}};
    \node[amn, fill=brainpink, text=pinkdark, visible on=<3->] at (3.0,0.3) {Անտերոգրադ\\ չես գրանցում \textbf{նորը}};
  \end{tikzpicture}
  \vspace{3mm}
  \begin{columns}[T]
    \begin{column}{0.47\textwidth}
      \onslide<2->{%
      \begin{card}[colback=lavender!35!cardwhite]
        \emoji{tv}\ Սերիալներում հերոսը գլխին հարված է ստանում ու մոռանում վերջին 20 տարին։ Իրականում սա \textbf{շատ հազվադեպ} է ու սովորաբար ժամանակավոր
      \end{card}}
    \end{column}
    \begin{column}{0.47\textwidth}
      \onslide<3->{%
      \begin{card}[colback=brainpink!25!cardwhite]
        \emoji{film}\ Ավելի հաճախ է, որովհետև հիպոկամպը փխրուն է։ Հ.Մ.-ը, «Memento» ֆիլմը։ Իսկ եթե «Memento»-ն դիտել ես, բայց չես հիշում, դա այնքան էլ չի օգնում (բայց հեգնական է)
      \end{card}}
    \end{column}
  \end{columns}
\end{frame}

\begin{frame}{Ինչն է վնասում հիպոկամպին}
  \begin{columns}[T]
    % Alzheimer's first, as in the book, with today's science (2026-11-12): the book has tangles only, no amyloid.
    % Same 2x2 grid as before; the science box and Burnett's quote share the bottom slot (\only), so nothing grows
    \begin{column}{0.47\textwidth}
      \begin{card}\emoji{brain}\ \textbf{Ալցհեյմեր}. \textbf{tau} սպիտակուցը խճճվում է այգու խողովակի պես\end{card}
      \onslide<3->{\begin{card}\emoji{hospital}\ \textbf{Ինսուլտ}. արյան կարճ ընդհատումն էլ բավական է\end{card}}
    \end{column}
    \begin{column}{0.47\textwidth}
      \onslide<3->{%
      \begin{card}\emoji{pill}\ \textbf{Հերպեսի վիրուսը} (շրթունքի «մրսածությունը»), հազվադեպ\end{card}
      \begin{card}\emoji{wine}\ \textbf{Կորսակովի համախտանիշ}. ալկոհոլիզմ, B1-ի պակաս\end{card}}
    \end{column}
  \end{columns}
  \only<1-3>{\onslide<2->{%
  \begin{scienceupdate}
    Նաև \textbf{ամիլոիդ}. 2023-ից այն մաքրող դեղերը մի քիչ օգնում են
  \end{scienceupdate}}}
  \only<4->{%
  \begin{card}[borderline west={1.4mm}{0pt}{accent}, fontupper=\mediumfont\small]
    \emoji{book}\ Բըրնեթ. «Մի պացիենտի ուղեղը վնասել էր քթի միջով խրված բիլիարդի փայտը։ Ոչ կոնտակտային սպորտ գոյություն չունի»։
  \end{card}}
  \source{van Dyck et al., NEJM, 2023; Sims et al., JAMA, 2023}
\end{frame}

\begin{frame}{Եվ այնուամենայնիվ...}
  \begin{columns}[c]
    \begin{column}{0.6\textwidth}
      \begin{itemize}
        \item Խճճված, փխրուն, անհետևողական, եսասեր
        \item Բայց \textbf{մեծ մասամբ աշխատում է}
        \item Տարողությամբ ու ճկունությամբ գերազանցում է ամենաժամանակակից սուպերհամակարգիչներին
      \end{itemize}
      \vspace{3mm}
      \begin{bookquote}
        Մարդու հիշողությունը կատարյալ չէ, բայց \textbf{բավականաչափ լավն է}։
      \end{bookquote}
    \end{column}
    \begin{column}{0.34\textwidth}
      \centering\mascot[3cm]{happy}
    \end{column}
  \end{columns}
\end{frame}
